\chapter{Alcance}
El presente anteproyecto tiene como alcance proponer la reactivación funcional
del robot cuadrúpedo Nova Spot Micro del laboratorio de robótica de
la Universidad Santo Tomás mediante la caracterización de su modelo cinemático y la implementación de un sistema electrónico de
control que permita la ejecución de patrones de locomoción cuadrúpeda.

En este sentido, se proyecta:



\begin{itemize}

\item Realizar la caracterización geométrica y estructural de la plataforma con el
fin de identificar los parámetros necesarios para su modelado.

\item Formular el modelo cinemático directo e inverso de las extremidades del
robot y de igual manera los jacobianos del sistema junto con el modelo dinámico como base para la generación de trayectorias articulares.

\item Integrar la arquitectura electrónica necesaria para el accionamiento y
sensado de las articulaciones.

\item Implementar una marcha nominal mediante una máquina de estados finitos,
orientada a la coordinación de movimientos en fases de apoyo y oscilación para
postura, paso y gateo.

\item Entrenar en simulación y evaluar tanto en simulación como en el robot
físico una política de aprendizaje por refuerzo que actúe como capa correctiva
sobre la marcha nominal.

\item Establecer un mecanismo HRI (human-robot-interface) para la operación del
robot.

\item Definir un esquema de validación experimental que permita evaluar la
estabilidad y funcionamiento en condiciones controladas de laboratorio.

\end{itemize}

El impacto esperado del proyecto radica en la habilitación operativa de la
plataforma como herramienta académica y de investigación, generando una base
técnica documentada que facilite futuros desarrollos en controles mas complejos y con
mejoras del desempeño en locomoción cuadrúpeda.

\section{Limitaciones del proyecto}


El presente proyecto se orienta a la reactivación funcional de la plataforma
cuadrúpeda mediante su caracterización mecánica, modelado cinemático y dinámico e
implementación de un sistema de control para los modos de locomoción tipo paso y gateo.
En coherencia con este enfoque, no se contempla el desarrollo de un sistema
completamente autónomo ni la integración de conectividad a internet,
comunicación en la nube. El control estará
basado en patrones nominales de marcha predefinidos y programados mediante una
máquina de estados finitos. Sobre estas referencias se evaluará una política de
aprendizaje por refuerzo encargada exclusivamente de generar correcciones
pequeñas y acotadas de estabilización. La política no comandará directamente
las señales PWM de los servomotores, no sustituirá los límites articulares ni
el supervisor de seguridad, y no incorporará navegación autónoma o toma de
decisiones en entornos dinámicos.

Asimismo, el proyecto no incluye el diseño ni la optimización avanzada del
sistema energético. No se realizará una caracterización detallada de baterías
orientada a maximizar densidad energética o autonomía prolongada, ya que la
selección del sistema de alimentación estará limitada por las restricciones de
peso y capacidad estructural de la plataforma existente.

En cuanto al desempeño locomotor, el robot estará orientado a operar en
superficies planas y condiciones controladas de laboratorio. No se contempla la
implementación de estrategias que permitan subir o bajar desniveles,
desplazarse en terrenos altamente irregulares o ejecutar maniobras complejas
en ambientes no estructurados.

El proyecto se orientará a locomoción de baja velocidad y no establecerá como
requisito una velocidad máxima de desplazamiento. La velocidad se limitará de
acuerdo con la estabilidad observada, la capacidad de los actuadores y los
criterios de seguridad definidos durante la caracterización. Se considerará
funcional el desplazamiento cuando el robot complete de forma verificable y
repetible los ciclos de paso y gateo establecidos; la velocidad alcanzada se
medirá y reportará como resultado experimental.

Finalmente, el proyecto se limitará al estudio y aplicación de modelos
cinemáticos para la generación de marchas estáticas tipo paso y gateo, sin
abordar un modelado dinámico complejo o muy robusto del sistema ni la implementación de
locomoción rápida como trote o galope. El entrenamiento de la política se
realizará primero en simulación y su transferencia al robot físico será
progresiva, con correcciones limitadas y una parada de emergencia independiente.

\section{Productos esperados}


\begin{enumerate}

\item Informe técnico final del proyecto, que documente de manera estructurada el proceso de caracterización mecánica, modelado cinemático, diseño electrónico e implementación del sistema de control de la plataforma cuadrúpeda.

\item Modelos cinemático y dinámico formales del robot, incluyendo el desarrollo matemático de la cinemática directa e inversa de las extremidades y de los jacobianos articulares, debidamente sustentados y validados en simulación o pruebas experimentales.

\item Implementación de la arquitectura electrónica de control, que incluya:

\begin{itemize}

\item Diagramas eléctricos.

\item Implementación de micro-computador necesario con sus respectivos componentes electrónicos para ejecutar los algoritmos de control.

\item Integración de actuadores y sensores necesarios para el control articular.

\end{itemize}

\item Código fuente documentado del sistema de control, correspondiente a la
generación y coordinación de postura, paso y gateo, así como al entrenamiento e
integración de la política correctiva de aprendizaje por refuerzo.

\item Prototipo funcional de la plataforma reactivada, capaz de mantener una
postura y ejecutar los patrones nominales de paso y gateo en condiciones
controladas de laboratorio, con evaluación comparativa de su desempeño con y
sin la capa correctiva aprendida.

\item Manual técnico de operación y puesta en marcha, que permita a futuros estudiantes o investigadores comprender la arquitectura del sistema y reutilizar la plataforma en proyectos posteriores.

\end{enumerate}
